% Options for packages loaded elsewhere
\PassOptionsToPackage{unicode}{hyperref}
\PassOptionsToPackage{hyphens}{url}
\documentclass[
]{article}
\usepackage{xcolor}
\usepackage[margin=1in]{geometry}
\usepackage{amsmath,amssymb}
\setcounter{secnumdepth}{-\maxdimen} % remove section numbering
\usepackage{iftex}
\ifPDFTeX
  \usepackage[T1]{fontenc}
  \usepackage[utf8]{inputenc}
  \usepackage{textcomp} % provide euro and other symbols
\else % if luatex or xetex
  \usepackage{unicode-math} % this also loads fontspec
  \defaultfontfeatures{Scale=MatchLowercase}
  \defaultfontfeatures[\rmfamily]{Ligatures=TeX,Scale=1}
\fi
\usepackage{lmodern}
\ifPDFTeX\else
  % xetex/luatex font selection
\fi
% Use upquote if available, for straight quotes in verbatim environments
\IfFileExists{upquote.sty}{\usepackage{upquote}}{}
\IfFileExists{microtype.sty}{% use microtype if available
  \usepackage[]{microtype}
  \UseMicrotypeSet[protrusion]{basicmath} % disable protrusion for tt fonts
}{}
\makeatletter
\@ifundefined{KOMAClassName}{% if non-KOMA class
  \IfFileExists{parskip.sty}{%
    \usepackage{parskip}
  }{% else
    \setlength{\parindent}{0pt}
    \setlength{\parskip}{6pt plus 2pt minus 1pt}}
}{% if KOMA class
  \KOMAoptions{parskip=half}}
\makeatother
\usepackage{color}
\usepackage{fancyvrb}
\newcommand{\VerbBar}{|}
\newcommand{\VERB}{\Verb[commandchars=\\\{\}]}
\DefineVerbatimEnvironment{Highlighting}{Verbatim}{commandchars=\\\{\}}
% Add ',fontsize=\small' for more characters per line
\usepackage{framed}
\definecolor{shadecolor}{RGB}{248,248,248}
\newenvironment{Shaded}{\begin{snugshade}}{\end{snugshade}}
\newcommand{\AlertTok}[1]{\textcolor[rgb]{0.94,0.16,0.16}{#1}}
\newcommand{\AnnotationTok}[1]{\textcolor[rgb]{0.56,0.35,0.01}{\textbf{\textit{#1}}}}
\newcommand{\AttributeTok}[1]{\textcolor[rgb]{0.13,0.29,0.53}{#1}}
\newcommand{\BaseNTok}[1]{\textcolor[rgb]{0.00,0.00,0.81}{#1}}
\newcommand{\BuiltInTok}[1]{#1}
\newcommand{\CharTok}[1]{\textcolor[rgb]{0.31,0.60,0.02}{#1}}
\newcommand{\CommentTok}[1]{\textcolor[rgb]{0.56,0.35,0.01}{\textit{#1}}}
\newcommand{\CommentVarTok}[1]{\textcolor[rgb]{0.56,0.35,0.01}{\textbf{\textit{#1}}}}
\newcommand{\ConstantTok}[1]{\textcolor[rgb]{0.56,0.35,0.01}{#1}}
\newcommand{\ControlFlowTok}[1]{\textcolor[rgb]{0.13,0.29,0.53}{\textbf{#1}}}
\newcommand{\DataTypeTok}[1]{\textcolor[rgb]{0.13,0.29,0.53}{#1}}
\newcommand{\DecValTok}[1]{\textcolor[rgb]{0.00,0.00,0.81}{#1}}
\newcommand{\DocumentationTok}[1]{\textcolor[rgb]{0.56,0.35,0.01}{\textbf{\textit{#1}}}}
\newcommand{\ErrorTok}[1]{\textcolor[rgb]{0.64,0.00,0.00}{\textbf{#1}}}
\newcommand{\ExtensionTok}[1]{#1}
\newcommand{\FloatTok}[1]{\textcolor[rgb]{0.00,0.00,0.81}{#1}}
\newcommand{\FunctionTok}[1]{\textcolor[rgb]{0.13,0.29,0.53}{\textbf{#1}}}
\newcommand{\ImportTok}[1]{#1}
\newcommand{\InformationTok}[1]{\textcolor[rgb]{0.56,0.35,0.01}{\textbf{\textit{#1}}}}
\newcommand{\KeywordTok}[1]{\textcolor[rgb]{0.13,0.29,0.53}{\textbf{#1}}}
\newcommand{\NormalTok}[1]{#1}
\newcommand{\OperatorTok}[1]{\textcolor[rgb]{0.81,0.36,0.00}{\textbf{#1}}}
\newcommand{\OtherTok}[1]{\textcolor[rgb]{0.56,0.35,0.01}{#1}}
\newcommand{\PreprocessorTok}[1]{\textcolor[rgb]{0.56,0.35,0.01}{\textit{#1}}}
\newcommand{\RegionMarkerTok}[1]{#1}
\newcommand{\SpecialCharTok}[1]{\textcolor[rgb]{0.81,0.36,0.00}{\textbf{#1}}}
\newcommand{\SpecialStringTok}[1]{\textcolor[rgb]{0.31,0.60,0.02}{#1}}
\newcommand{\StringTok}[1]{\textcolor[rgb]{0.31,0.60,0.02}{#1}}
\newcommand{\VariableTok}[1]{\textcolor[rgb]{0.00,0.00,0.00}{#1}}
\newcommand{\VerbatimStringTok}[1]{\textcolor[rgb]{0.31,0.60,0.02}{#1}}
\newcommand{\WarningTok}[1]{\textcolor[rgb]{0.56,0.35,0.01}{\textbf{\textit{#1}}}}
\usepackage{graphicx}
\makeatletter
\newsavebox\pandoc@box
\newcommand*\pandocbounded[1]{% scales image to fit in text height/width
  \sbox\pandoc@box{#1}%
  \Gscale@div\@tempa{\textheight}{\dimexpr\ht\pandoc@box+\dp\pandoc@box\relax}%
  \Gscale@div\@tempb{\linewidth}{\wd\pandoc@box}%
  \ifdim\@tempb\p@<\@tempa\p@\let\@tempa\@tempb\fi% select the smaller of both
  \ifdim\@tempa\p@<\p@\scalebox{\@tempa}{\usebox\pandoc@box}%
  \else\usebox{\pandoc@box}%
  \fi%
}
% Set default figure placement to htbp
\def\fps@figure{htbp}
\makeatother
\setlength{\emergencystretch}{3em} % prevent overfull lines
\providecommand{\tightlist}{%
  \setlength{\itemsep}{0pt}\setlength{\parskip}{0pt}}
\usepackage{bookmark}
\IfFileExists{xurl.sty}{\usepackage{xurl}}{} % add URL line breaks if available
\urlstyle{same}
\hypersetup{
  pdftitle={Differential Expression with edgeR},
  hidelinks,
  pdfcreator={LaTeX via pandoc}}

\title{Differential Expression with edgeR}
\author{}
\date{\vspace{-2.5em}2026-04-17}

\begin{document}
\maketitle

\subsection{Introduction}\label{introduction}

edgeR is one of the three dominant RNA-seq differential-expression
frameworks (alongside DESeq2 and limma-voom). It fits a
negative-binomial generalised linear model to each gene's counts, with
empirical-Bayes shrinkage borrowing information across genes to
stabilise per-gene dispersion estimates. The combination of NB
likelihood and shrinkage gives edgeR its characteristic strength on
small samples and its competitive accuracy on standard RNA-seq
experiments.

\subsection{Prerequisites}\label{prerequisites}

A working understanding of negative-binomial regression, RNA-seq
normalisation (TMM), generalised linear models with offsets, and the
empirical-Bayes shrinkage approach.

\subsection{Theory}\label{theory}

Per gene \(i\) and sample \(j\):

\[Y_{ij} \sim \mathrm{NB}(\mu_{ij}, \phi_i), \qquad \log \mu_{ij} = \log S_j + x_j^\top \beta_i,\]

with \(S_j\) the TMM-normalised library size offset and \(\phi_i\) the
gene-wise dispersion. Empirical-Bayes shrinkage (\texttt{estimateDisp})
pulls per-gene dispersions toward a common trend across genes,
stabilising estimation.

Three test choices:

\begin{itemize}
\tightlist
\item
  \textbf{exactTest}: NB-exact test for two-group comparisons; rarely
  used today.
\item
  \textbf{glmQLFTest}: quasi-likelihood F-test; the recommended default
  for nearly all designs because it accounts for uncertainty in
  dispersion estimation.
\item
  \textbf{glmLRT}: classical likelihood-ratio test; more aggressive but
  anti-conservative when sample sizes are very small.
\end{itemize}

\subsection{Assumptions}\label{assumptions}

Counts approximately negative-binomial conditional on covariates; the
design matrix encodes the biological and technical factors correctly;
sufficient sample size for stable dispersion estimation (typically
\(\ge 3\) per group).

\subsection{R Implementation}\label{r-implementation}

\begin{Shaded}
\begin{Highlighting}[]
\FunctionTok{library}\NormalTok{(edgeR)}

\FunctionTok{set.seed}\NormalTok{(}\DecValTok{2026}\NormalTok{)}
\NormalTok{counts }\OtherTok{\textless{}{-}} \FunctionTok{matrix}\NormalTok{(}\FunctionTok{rpois}\NormalTok{(}\DecValTok{20000} \SpecialCharTok{*} \DecValTok{12}\NormalTok{, }\AttributeTok{lambda =} \DecValTok{30}\NormalTok{), }\DecValTok{20000}\NormalTok{, }\DecValTok{12}\NormalTok{)}
\NormalTok{group }\OtherTok{\textless{}{-}} \FunctionTok{factor}\NormalTok{(}\FunctionTok{rep}\NormalTok{(}\FunctionTok{c}\NormalTok{(}\StringTok{"A"}\NormalTok{, }\StringTok{"B"}\NormalTok{), }\AttributeTok{each =} \DecValTok{6}\NormalTok{))}

\NormalTok{y }\OtherTok{\textless{}{-}} \FunctionTok{DGEList}\NormalTok{(}\AttributeTok{counts =}\NormalTok{ counts, }\AttributeTok{group =}\NormalTok{ group)}
\NormalTok{y }\OtherTok{\textless{}{-}} \FunctionTok{calcNormFactors}\NormalTok{(y)}

\NormalTok{design }\OtherTok{\textless{}{-}} \FunctionTok{model.matrix}\NormalTok{(}\SpecialCharTok{\textasciitilde{}}\NormalTok{ group)}
\NormalTok{y }\OtherTok{\textless{}{-}} \FunctionTok{estimateDisp}\NormalTok{(y, design)}

\NormalTok{fit }\OtherTok{\textless{}{-}} \FunctionTok{glmQLFit}\NormalTok{(y, design)}
\NormalTok{qlf }\OtherTok{\textless{}{-}} \FunctionTok{glmQLFTest}\NormalTok{(fit, }\AttributeTok{coef =} \DecValTok{2}\NormalTok{)}

\FunctionTok{topTags}\NormalTok{(qlf)}
\FunctionTok{summary}\NormalTok{(}\FunctionTok{decideTests}\NormalTok{(qlf))}
\end{Highlighting}
\end{Shaded}

\subsection{Output \& Results}\label{output-results}

\texttt{topTags()} returns the top differentially-expressed genes
ordered by FDR-adjusted \(p\)-value, with \(\log_2\) fold change,
log-CPM, \(F\)-statistic, and BH-adjusted \(p\)-value.
\texttt{decideTests()} summarises significance calls per direction (up /
down / not significant) for cross-method comparison.

\subsection{Interpretation}\label{interpretation}

A reporting sentence: ``edgeR's quasi-likelihood F-test identified 342
differentially expressed genes at FDR \textless{} 0.05 between
conditions A and B (n = 6 per group), with 180 upregulated and 162
downregulated in condition B; the median absolute \(\log_2\) fold change
among significant genes was 1.4. Filtering with \texttt{filterByExpr}
retained 14\{,\}500 of 20\{,\}000 genes for testing.'' Always state the
FDR threshold, sample sizes, and pre-filtering decisions.

\subsection{Practical Tips}\label{practical-tips}

\begin{itemize}
\tightlist
\item
  Use \texttt{glmQLFTest()} as the default for nearly all designs; it
  has better Type I error control than \texttt{glmLRT()} at small sample
  sizes.
\item
  Filter low-count genes with \texttt{filterByExpr()} before dispersion
  estimation; testing 20\{,\}000 mostly-zero genes wastes power and
  produces unstable shrinkage.
\item
  For paired or blocked designs, include the blocking factor in the
  \texttt{design} matrix; the QLF test handles the corresponding
  contrast correctly.
\item
  Use \texttt{glmTreat()} for fold-change-thresholded inference (testing
  \textbar LFC\textbar{} \textgreater{} a chosen cutoff rather than LFC
  ≠ 0); this is more biologically meaningful than testing strict zero.
\item
  Visualise with MA plot, volcano plot, and a heatmap of top genes;
  \texttt{Glimma::glimmaMA()} provides interactive versions for
  exploration.
\item
  Compare edgeR results with DESeq2 and limma-voom on the same data;
  concordance across methods is reassuring, while disagreement signals
  model-sensitive genes that warrant investigation.
\end{itemize}

\subsection{R Packages Used}\label{r-packages-used}

\texttt{edgeR} for \texttt{DGEList}, \texttt{calcNormFactors},
\texttt{estimateDisp}, \texttt{glmQLFit}, \texttt{glmQLFTest}, and the
canonical NB workflow; \texttt{limma} for the voom alternative;
\texttt{DESeq2} for the third major DE method; \texttt{Glimma} for
interactive MA and volcano plots; \texttt{EnhancedVolcano} for
publication-ready volcano plots.

\subsection{For Reviewers}\label{for-reviewers}

\textbf{What to look for in a paper using this method.}

\begin{itemize}
\tightlist
\item
  \textbf{Common misapplications.}
\item
  \textbf{Diagnostics that should be reported but often aren't.}
\item
  \textbf{Red flags in tables and figures.}
\item
  \textbf{What to verify.}
\item
  \textbf{What an adequate Methods paragraph must contain.}
\end{itemize}

\subsection{See also --- labs in this
chapter}\label{see-also-labs-in-this-chapter}

\begin{itemize}
\tightlist
\item
  \href{labs/lab-rna-seq-with-deseq2-edger.Rmd}{RNA-seq with DESeq2 /
  edgeR}
\item
  \href{labs/lab-enrichment-analysis-gsea-ora.Rmd}{Enrichment analysis
  (GSEA / ORA)}
\item
  \href{labs/lab-scrna-seq-with-seurat.Rmd}{scRNA-seq with Seurat}
\end{itemize}

\end{document}
